\section{Energy bounds and strong locality}\label{sec:energy}

Our first analytic task is to construct the vacuum net.  We do this
without assuming that any module has already been integrated to a net
representation.  The argument has two stages.  A finite set of
quasi-primary generators first gives polynomial energy bounds of
unspecified order.  A second estimate improves the order to the conformal
weight minus one, which is the order needed for the strong-locality
criterion.

\begin{theorem}\label{thm:energy}
Let $V$ be a simple unitary strongly rational vertex operator algebra,
and let $W$ be a grading-restricted unitary $V$-module.  Every vertex
operator on $W$ satisfies polynomial energy bounds.  More precisely, if
$a\in V_d$ is quasi-primary, $d\geq2$, and
$(-1)^d\Theta a=a$, then there are constants $C_a>0$ and
$p_a\in\mathbb Z_{\geq0}$ such that
\begin{equation}\label{eq:optimal-vertex-bound}
 \|a[m]\xi\|
 \leq C_a(1+|m|)^{p_a}
       \|(1+L_0)^{d-1}\xi\|,
 \qquad m\in\mathbb Z,\quad \xi\in W.
\end{equation}
Weight-one fields and the Virasoro field satisfy energy bounds of at
most linear order.  The constants may depend on $V$, $W$, and $a$.
\end{theorem}

The scalar VOA satisfies the theorem directly.  We assume for the proof
that $V\ne\mathbb C\Omega$, and fix $W$.  Since $W$ is a finite
orthogonal sum of ordinary unitary modules, it suffices to establish all
estimates on those summands and take the maximum of the finitely many
constants.  Write $W(h)$ for the energy-$h$ subspace, with $W(h)=0$ when
$h$ is absent.  These spaces are finite dimensional and mutually
orthogonal.  Unitarity of the positive-energy Virasoro representation
gives $h\geq0$.  We use the spectral projections
\[
 P_N=\mathbf1_{[0,N]}(L_0),\qquad
 \Pi_h=\mathbf1_{\{h\}}(L_0),\qquad H=1+L_0.
\]
All mode identities below initially act on finite-energy vectors.
For an Hermitian quasi-primary $a$ of weight $d$, the invariant form
gives
\begin{equation}\label{eq:energy-adjoint}
 a[m]^*=a[-m],\qquad
 [L_j,a[m]]=((d-1)j-m)a[j+m]\quad(j=-1,0,1),
\end{equation}
where the adjoint is taken between the indicated energy spaces.

\subsection{A finite generating filtration}

Choose an integer $D\geq2$ such that
$V_w=C_2(V)_w$ for all $w>D$.  In each weight $1,\ldots,D$, choose a
finite set of Hermitian vectors spanning the quasi-primary space, and
let $S$ be their union.  Such a choice is possible because
$(-1)^w\Theta$ is an antilinear involution of that space.  We distinguish
the conformal weight of a monomial from the sum of the weights of its
generators.

\begin{lemma}\label{lem:energy-filtration}
Every homogeneous vector of weight $w$ is a linear combination of
monomials
\begin{equation}\label{eq:energy-monomial}
 b=a^1_{(-1-n_1)}\cdots a^\ell_{(-1-n_\ell)}\Omega,
 \qquad a^i\in S,\quad n_i\geq0,
\end{equation}
with
\[
 D_b:=\sum_{i=1}^{\ell}d_{a^i},\qquad
 w=D_b+\sum_{i=1}^{\ell}n_i.
\]
If the vector belongs to $C_2(V)_w$, the monomials can be chosen with
$D_b\leq w-1$.
\end{lemma}

\begin{proof}
Let $E_t$ be the span of arbitrary-mode strings of elements of $S$ on
$\Omega$ whose total generator weight is at most $t$.  Induction on
string length, using the iterate identity, gives
\begin{equation}\label{eq:energy-filtration-stability}
 u_{(r)}E_s\subset E_{t+s}\qquad(u\in E_t, r\in\mathbb Z).
\end{equation}
For the induction step, write a chosen string as $u=a_{(m)}u'$, with
$a\in S$ and $u'$ of generator degree at most $t-d_a$.
The iterate formula expresses $u_{(r)}v$, $v\in E_s$, as a sum of terms
of the two forms
\[
 a_{(m-j)}u'_{(r+j)}v,
 \qquad u'_{(m+r-j)}a_{(j)}v,
 \qquad j\geq0.
\]
For the first, apply induction to $u'$ and then prepend the generator
mode. For the second, $a_{(j)}v\in E_{s+d_a}$, so induction on $u'$
again gives a vector in $E_{t+s}$. Each of the two sums truncates on
the fixed vector $v$.

We first prove $V_w\subset E_w$ by induction on $w$.  At weights at
most $D$, use the orthogonal decomposition
\[
 V_w=(V_w\cap\ker L_1)+L_{-1}V_{w-1}.
\]
The first summand is spanned by our chosen generators, and differentiation
does not increase generator weight.  A homogeneous term $u_{(-2)}v$ of
weight $w$ has $\mathrm{wt}(u)+\mathrm{wt}(v)=w-1$.  If both factors
have positive weight, induction and
\eqref{eq:energy-filtration-stability} put it in $E_{w-1}$.  If
$v=\Omega$, it is $L_{-1}u$ and has the same property; terms with
$u\in\mathbb C\Omega$ vanish.  This proves
$C_2(V)_w\subset E_{w-1}$ and completes the induction above $D$ as well.

Now order an arbitrary-mode string by moving nonnegative ordinary modes
to the right.  A commutator of modes of generators $a,b$ involves
$a_{(j)}b$, $j\geq0$, of weight at most $d_a+d_b-1$.  The inclusion
just proved places this state in $E_{d_a+d_b-1}$.  Thus every
commutator term lowers the total generator weight, whereas the
uncommuted term moves the chosen nonnegative mode one place to the
right.  Induction on total generator weight, and then on the number of
such moves, puts the string into negative ordinary-mode form.  A
nonnegative mode reaching $\Omega$ vanishes.  Taking homogeneous
components proves both assertions.
\end{proof}

We shall use only finitely many expansions from
Lemma~\ref{lem:energy-filtration} at a time.  Constants may depend on
these fixed expansions.  In particular, no uniformity over all
monomials is asserted.

Here is the mode expansion that controls the intermediate energies.
Iterating the negative-product formula from the outside of
\eqref{eq:energy-monomial} places each generator either on the left,
with physical index
\begin{equation}\label{eq:energy-left-right-indices}
 -d_{a^i}-n_i-j_i,
 \qquad\text{or on the right, with index}\qquad
 1-d_{a^i}+j_i,\qquad j_i\geq0.
\end{equation}
The coefficients are products of signed binomial coefficients, bounded
by $C_b\prod_i(1+j_i)^{n_i}$.  There are finitely many order patterns;
in each pattern every left factor precedes every right factor.  The
vacuum mode imposes the constraint that the physical indices sum to the
desired mode of $b$.

In a matrix element with both endpoint energies at most $N$, move the
left factors to the testing vector by adjunction.  They then all lower
energy.  Each right factor can raise energy by at most $d_{a^i}-1$.
Every intermediate vector therefore lies in $P_{N+c_b}W$ for a fixed
integer $c_b$, and every relevant $j_i$ is $O_b(1+N)$.  This proves
finiteness of the sums used below and permits the insertion of those
fixed intermediate cutoffs.

\subsection{Polynomial energy bounds}

For integer $N\geq1$, define
\[
 B_N=\max\left(1,
       \max_{a\in S,\,m\in\mathbb Z}
          \|P_Na[m]P_N\|^{1/d_a}\right).
\]
The $d_a$th root measures growth per unit of generator weight, so a
monomial of generator degree $D_b$ contributes the power
$B_{N+c_b}^{D_b}$.
Only $|m|\leq N$ can contribute.  The preceding mode expansion gives,
for each fixed homogeneous $u\in V_w$,
\begin{equation}\label{eq:rough-monomial-bound}
 \|P_Nu[\mu]P_N\|
 \leq C_u(1+N)^{p_u}B_{N+c_u}^{w}.
\end{equation}
All cutoffs and exponents in this preliminary estimate are fixed once
the expression for $u$ is fixed.

We need an initial growth bound before improving
\eqref{eq:rough-monomial-bound}.  For $a\in S$ of weight $d$, set
$T_a=(a_{(-1)}a)[0]$.  The normal-product formula reads
\begin{equation}\label{eq:positive-normal-product}
 T_a=\sum_{j\geq0}
 \bigl(a[-d-j]a[d+j]+a[d-1-j]a[j+1-d]\bigr).
\end{equation}
Every summand is a positive square as a finite-energy quadratic form.
The second family contains $a[m]^*a[m]$ for every $m\geq0$.
Consequently
\begin{equation}\label{eq:positive-column}
 \sum_{m\geq0}\|a[m]\xi\|^2
 \leq \langle T_a\xi,\xi\rangle.
\end{equation}
Zhu's ordinary-module one-point convergence theorem applies to
$a_{(-1)}a$ \cite[Sections~4.3--4.4 and Theorem~4.4.1]{Zhu}.  Its finiteness
Condition~C holds here: besides $C_2$-cofiniteness and finite-dimensional
grades, the required spanning by Virasoro descendants of primary
vectors follows from unitarity.  For fixed $0<q<1$ we therefore have
\[
 Z_a(q):=\operatorname{Tr}_W(T_aq^{L_0})<\infty.
\]
Here the trace is the convergent graded sum
$\sum_h q^h\operatorname{Tr}_{W(h)}(T_a|_{W(h)})$; no bounded
operator extension of $T_a$ is presumed.
Positivity on each energy level and \eqref{eq:positive-column} give
\[
 \sup_{\substack{\xi\in W(h)\\\|\xi\|=1}}
       \sum_{m\geq0}\|a[m]\xi\|^2
 \leq q^{-h}Z_a(q).
\]
Negative compressed modes are controlled by adjunction.  It follows
that $\log B_N=O(N+1)$.  This is the positive-trace method used for the
vacuum exponential bound in
\cite[Lemma~3.5 and Proposition~3.17]{CarpiTomassini}; the ordinary-module
trace theorem gives the version needed here.

For $a\in S$ of weight $d$ and $m>0$, the mode commutator is
\begin{equation}\label{eq:self-mode-commutator}
 C_{a,m}:=[a[m],a[-m]]
   =\sum_{j\geq0}\binom{d-1+m}{j}(a_{(j)}a)[0].
\end{equation}
The nonzero states on the right have weights at most $2d-1$.
Thus \eqref{eq:rough-monomial-bound}, uniformly for $1\leq m\leq N$,
implies
\begin{equation}\label{eq:rough-commutator}
 \|C_{a,m}P_N\|
 \leq C(1+N)^p B_{N+c}^{2d-1}.
\end{equation}

The descent from a commutator estimate to a mode estimate is the
mechanism used in \cite[Proposition~3.3]{CTW}. We record it explicitly
with its input-energy cutoff, which enters the subsequent iteration. Put
$A_h=\|a[-m]:W(h)\longrightarrow W(h+m)\|$, interpreting absent
levels as zero.  Adjunction and \eqref{eq:self-mode-commutator} give
\begin{equation}\label{eq:mode-norm-descent}
 A_h^2\leq A_{h-m}^2+\|C_{a,m}|_{W(h)}\|.
\end{equation}
Indeed, apply the identity to a unit vector in $W(h)$, and note that
$\|a[m]|_{W(h)}\|=A_{h-m}$.  Iterating descends to negative energy in
at most $1+N/m$ steps when $h\leq N$.  There is no factor depending on
the dimensions of the levels.  The raising mode may have output above
$N$; only the input level enters the commutator estimate.

The same argument applied to
$[L_m,L_{-m}]=2mL_0+\frac c{12}(m^3-m)$ gives
\begin{equation}\label{eq:virasoro-linear-energy}
 \|L_m\xi\|\leq C(1+|m|)^2\|H\xi\|.
\end{equation}
In particular $L_{\pm1}$ have linear energy bounds.  For the zero mode
of $a$, use
\begin{equation}\label{eq:zero-mode-rough-transfer}
 d\,a[0]=[L_1,a[-1]].
\end{equation}
The two products involve only a fixed shift of input energy.  Combining
\eqref{eq:rough-commutator}--\eqref{eq:zero-mode-rough-transfer}, and
enlarging a common fixed cutoff, yields
\begin{equation}\label{eq:first-energy-contraction}
 B_N\leq C(1+N)^p B_{N+c}^{\rho},
 \qquad \rho=1-\frac1{2D}<1.
\end{equation}

The use of an exponential starting bound to close a contraction with an
additive cutoff also appears in the proof of
\cite[Theorem~4.5]{CarpiTomassini}. Here the finite-$C_2$ filtration
supplies the contracting estimate. The contraction remains effective
despite the increasing cutoff: after $r$ iterations,
\[
 \log B_N\leq
 \sum_{j=0}^{r-1}\rho^j
       \bigl(\log C+p\log(1+N+jc)\bigr)
       +\rho^r\log B_{N+rc}.
\]
The final term tends to zero by the exponential starting bound.  The
inequality $1+N+jc\leq(1+N)(1+jc)$ bounds the sum by
$C'+\frac p{1-\rho}\log(1+N)$.  Hence $B_N$ grows polynomially.

For an arbitrary mode $u[\mu]$ on $W(h)$, apply
\eqref{eq:rough-monomial-bound} with an endpoint cutoff of size
$O(1+h+|\mu|)$.  This bounds its norm by a polynomial in
$1+h+|\mu|$.  Distinct input energy levels have distinct output levels
under a fixed mode, so the squared norms add.  We have proved polynomial
energy bounds for every vertex operator on $W$.

In particular, smooth smears and their formal adjoints map
$\mathcal H_W^\infty=\bigcap_{r\geq0}\mathcal D(H^r)$ continuously
into itself: the mode shifts cost only polynomial factors, while smooth
Fourier coefficients decrease faster than any power.  Spectral
truncation shows that closing a smear from $W$ or from
$\mathcal H_W^\infty$ gives the same operator.  These facts will justify
the covariance calculation below.

\subsection{A sharper normal-product estimate}

Normal-product estimates also prove the closure of polynomial energy
bounds under homogeneous generation \cite[Proposition~6.1]{CKLW}.
The estimate below keeps the specific energy power needed for the
second contraction and locality. We divide the zero-mode norm by its
target power $N^{d-1}$ for $d\geq2$. The column normalization is
chosen so that its square has the desired energy power $N^{2d-1}$
for the positive zero mode $T_a$. Define
\[
 R_a(N)=\sup_{\substack{\xi\in P_NW\\\|\xi\|\leq1}}
          \left(\sum_{m\geq0}\|a[m]\xi\|^2\right)^{1/2}
\]
and, for integer $N\geq1$,
\begin{equation}\label{eq:sharp-energy-control}
 X_N=\max\left(
 1,
 \max_{a\in S}\left(\frac{R_a(N)}{N^{d_a-1/2}}\right)^{1/d_a},
 \max_{\substack{a\in S\\d_a\geq2}}
       \left(\frac{\|a[0]P_N\|}{N^{d_a-1}}\right)^{1/d_a}
 \right).
\end{equation}
The preceding bounds show that $X_N$ grows at most polynomially.
The column contains at most $1+N$ nonzero modes.

\begin{lemma}\label{lem:sharp-normal-product}
For a monomial $b$ in \eqref{eq:energy-monomial} of weight $w\geq2$,
there are constants $C_b,c_b$ such that
\begin{equation}\label{eq:sharp-normal-product}
 \|b[0]P_N\|
 \leq C_b N^{w-1}X_{N+c_b}^{D_b}\qquad(N\geq1).
\end{equation}
At weights zero and one the corresponding bounds are a constant and
$C_bN^{1/2}X_N$, respectively.
\end{lemma}

\begin{proof}
Use the left/right expansion \eqref{eq:energy-left-right-indices} for
the zero mode, between unit vectors $\eta,\xi\in P_NW$.  Consider first
an order pattern having $p\geq1$ left and $q\geq1$ right factors.  Let
$D_L,D_R$ be the sums of generator weights on those sides, and let
$L(\boldsymbol j_L),R(\boldsymbol j_R)$ denote their ordered products
of modes.  Write
\[
 A_{\boldsymbol j_L}=\|L(\boldsymbol j_L)^*\eta\|,
 \qquad
 B_{\boldsymbol j_R}=\|R(\boldsymbol j_R)\xi\|.
\]
Repeated use of the column bound defining $X$, with the fixed
intermediate cutoff already established, gives
\begin{align}
 \sum_{\boldsymbol j_L}A_{\boldsymbol j_L}^2
  &\leq C N^{2D_L-p}X_{N+c_b}^{2D_L},\label{eq:left-path-squares}\\
 \sum_{\boldsymbol j_R}B_{\boldsymbol j_R}^2
  &\leq C N^{2D_R-q}X_{N+c_b}^{2D_R}.
       \label{eq:right-path-squares}
\end{align}
To include the finitely many negative right indices, observe that for a
fixed positive integer $s$,
\[
 \|a[-s]P_M\|=\|P_Ma[s]\|
       \leq\|a[s]P_{M+s}\|\leq R_a(M+s).
\]
Thus adding those finitely many indices to the column costs only a
constant and a fixed cutoff enlargement.  Equations
\eqref{eq:left-path-squares}--\eqref{eq:right-path-squares} follow by
applying this one-factor estimate successively to the vectors on each
path and then summing.  Orthogonality between different paths is not
needed.

Let $\mathcal L,\mathcal R$ be the sets of left and right factors.
By \eqref{eq:energy-left-right-indices}, the zero-mode constraint is
\[
 \sum_{i\in\mathcal R}j_i-\sum_{i\in\mathcal L}j_i
   =D_b+\sum_{i\in\mathcal L}n_i-q.
\]
Thus one index is determined by the others. Given
a left tuple, at most $CN^{q-1}$ right tuples complete it: the last
right index is determined by the other $q-1$ indices.  Conversely,
there are at most $CN^{p-1}$ left completions of a right tuple.
Cauchy--Schwarz over the constrained pairs therefore gives
\begin{align*}
 \sum_{\text{constraint}}
       A_{\boldsymbol j_L}B_{\boldsymbol j_R}
 &\leq
 \left(CN^{q-1}\sum_{\boldsymbol j_L}A_{\boldsymbol j_L}^2\right)^{1/2}
 \left(CN^{p-1}\sum_{\boldsymbol j_R}B_{\boldsymbol j_R}^2\right)^{1/2}\\
 &\leq C N^{D_L+D_R-1}X_{N+c_b}^{D_L+D_R}.
\end{align*}
The constraint has saved one power of $N$ from the unconstrained count.
The binomial coefficients cost at most $N^{\sum_i n_i}$, producing
exactly the power $w-1$ in \eqref{eq:sharp-normal-product}.

An all-left pattern cannot have total physical index zero.  In an
all-right pattern, the indices have fixed lower bounds and sum to
zero, so every index and the number of terms are bounded independently
of $N$.  If $\ell\geq2$, individual column bounds give a power
$D_b-\ell/2\leq D_b-1$, which suffices.  If $\ell=1$ and $n_1=0$,
use the zero-mode term in $X_N$.  If $\ell=1$ and $n_1>0$, the weaker
power $d_{a^1}-1/2$ is bounded by $w-1$; differentiation at the zero
physical mode has a fixed scalar coefficient.  The vacuum and
weight-one cases follow directly from their defining bounds.
\end{proof}

For $a\in S$ of weight $d\geq2$, subtract the identity
\eqref{eq:self-mode-commutator} at $m=0$ from that at $m>0$.  The
former is the zero operator.  This removes the $j=0$ term without
requiring the state $a_{(0)}a$ to vanish.  Only fields of weights at
most $2d-2$ remain, with coefficients polynomial in $m$.
Lemma~\ref{lem:sharp-normal-product} now gives
\begin{equation}\label{eq:sharp-commutator}
 \|C_{a,m}P_N\|
 \leq C(1+m)^p N^{2d-3}X_{N+c}^{2d-2}.
\end{equation}
The weight-one estimate also fits this bound because
$N^{1/2}\leq N^{2d-3}$; the weight-zero term is bounded.  For all
other terms the generator degree is at most their conformal weight.
The descent \eqref{eq:mode-norm-descent}, with at most $1+N/m$
steps, consequently proves
\begin{equation}\label{eq:sharp-nonzero-mode}
 \|a[m]P_N\|
 \leq C(1+|m|)^pN^{d-1}X_{N+c}^{d-1},
 \qquad m\ne0.
\end{equation}
For a negative mode this follows directly by applying the raising-mode
recurrence at input energy at most $N$.  Thus the cutoff in $X$ does
not grow with $|m|$.  Positive modes of index greater than $N$ vanish.

\subsection{Transferring the bound to the zero mode}

The commutator \eqref{eq:zero-mode-rough-transfer} would lose one energy
power.  Möbius covariance transfers the estimate without that loss.
For Hermitian primary fields, \cite[Proposition~3.1]{CTW} transfers
an energy bound from one Fourier mode to all smooth smearings without
increasing its energy order. We give the quantitative cutoff transfer
for quasi-primary generators because its uniformity is needed in the
last iteration.

Let $a\in S$ have weight $d\geq2$.  Positive-energy unitary Möbius
integration acts on $\mathcal H_W$; it does not require a net
representation.  With the physical smearing convention, the test
$f(e^{i\vartheta})=e^{i\vartheta}$ selects $a[1]$.  Write
$a(f)=Y^W(a,f)$ for this module smear.  The
quasi-primary commutators \eqref{eq:energy-adjoint} integrate on the
smooth domain to
\[
 U_g a(f)U_g^*=a(f_g),\qquad
 f_g(g(\vartheta))=g'(\vartheta)^{d-1}f(e^{i\vartheta}).
\]
One obtains this identity by differentiating the conjugation with the
inversely evolving test; the preliminary polynomial bounds justify the
differentiation.  The positive-energy Möbius representation and these
covariance facts are recalled in \cite[Section~2]{CTW}.

Take the flow generated by $(L_1-L_{-1})/2$.  Its infinitesimal action
on $a[1]$ is
\[
 \left[\frac{L_1-L_{-1}}2,a[1]\right]
       =\frac{d-2}{2}a[2]+\frac d2a[0].
\]
The mean of $f_g$ therefore has nonzero derivative at the identity.
Fix a small $g$ for which that mean, denoted $c_g$, is nonzero.
Compression to one energy level gives
\begin{equation}\label{eq:mobius-zero-compression}
 \Pi_hU_ga[1]U_g^*\Pi_h=c_g a[0]\Pi_h.
\end{equation}

The conjugate $H'=U_gHU_g^*$ is a fixed linear combination of
$1,L_0,L_1,L_{-1}$.  By \eqref{eq:virasoro-linear-energy},
\[
 H'P_MW\subset P_{M+1}W,\qquad
 \|H'P_M\|\leq C_g(1+M).
\]
Consequently, for every integer $j\geq1$,
\begin{equation}\label{eq:mobius-growing-moment}
 \|H^jU_g^*P_N\|
 =\|(H')^jP_N\|
 \leq C_g^j\prod_{i=0}^{j-1}(1+N+i)
 \leq C_g^j(1+N+j)^j.
\end{equation}
For fixed powers, the same product estimate gives the usual Sobolev
continuity of $U_g$ and $U_g^*$, and hence preservation of
$\mathcal H_W^\infty$.  For a growing power it controls a spectral
tail.  Choose an integer $K>4C_g$.  If $r$ is fixed and $N>r$, then
\begin{align}
 \|H^r(1-P_{KN})U_g^*P_N\|
 &\leq (KN)^{-(N-r)}\|H^NU_g^*P_N\|\notag\\
 &\leq C_rN^r(3C_g/K)^N.
       \label{eq:mobius-tail}
\end{align}
The finitely many remaining $N$ are absorbed into the constant.  In
particular, the tail is bounded in every fixed graph norm, with $K$
independent of $N$ and of the dimension of $P_NW$.

Choose $r$ at least the preliminary polynomial energy order of $a[1]$.
Equations \eqref{eq:sharp-nonzero-mode} and \eqref{eq:mobius-tail}
give
\begin{align*}
 \|a[1]U_g^*P_N\|
 &\leq \|a[1]P_{KN}\|
      +\|a[1]H^{-r}\|\,
          \|H^r(1-P_{KN})U_g^*P_N\|\\
 &\leq C N^{d-1}X_{KN+c}^{d-1}+O(1).
\end{align*}
Compress using \eqref{eq:mobius-zero-compression}.  Since $a[0]$
preserves energy, taking the supremum of its level norms proves
\begin{equation}\label{eq:sharp-zero-mode}
 \|a[0]P_N\|
 \leq C N^{d-1}X_{KN+c}^{d-1}.
\end{equation}
The bounded error is absorbed because $X\geq1$ and $N\geq1$.

\subsection{Closing the optimal estimate}

Return to the positive operator $T_a$ in
\eqref{eq:positive-normal-product}, for any generator of weight $d$,
including $d=1$.  If $2d>D$, its defining state lies in
$C_2(V)_{2d}$, so Lemma~\ref{lem:energy-filtration} supplies an
expansion with generator degree at most $2d-1$.  If $2d\leq D$,
decompose that state into a quasi-primary part and an element of
$L_{-1}V_{2d-1}$.  The latter is in $C_2(V)$ and has the same strict
degree saving.  The former is a fixed combination of generators of
weight $2d$ and is controlled by \eqref{eq:sharp-zero-mode}.  In both
cases,
\begin{equation}\label{eq:sharp-positive-product}
 \|T_aP_N\|
 \leq C N^{2d-1}X_{K_0N}^{2d-1}
\end{equation}
for a sufficiently large fixed integer $K_0$.

Here and below the different enlarged cutoffs can indeed be replaced
by a common one.  Although $X_N$ itself need not be monotone, all its
unnormalized column and zero-mode numerators are monotone.  If
$N\leq M\leq cN$, its definition gives $X_N\leq cX_M$.
The finitely many fixed cutoff enlargements therefore change only the
constant in \eqref{eq:sharp-positive-product}.

Take the square root in \eqref{eq:positive-column} and
\eqref{eq:sharp-positive-product}, and then the $d$th root in the
definition of $X_N$.  The column term contracts with exponent
$(2d-1)/(2d)$.  Equation~\eqref{eq:sharp-zero-mode} contracts each
zero-mode term with exponent $(d-1)/d$.  Thus
\begin{equation}\label{eq:second-energy-contraction}
 X_N\leq C X_{K_0N}^{\sigma},\qquad
 \sigma=1-\frac1{2D}<1.
\end{equation}
The first contraction allowed an additive cutoff enlargement and used
an exponential starting bound. Here the cutoff grows multiplicatively;
the polynomial bound already proved for $X_N$ suffices. Namely,
\[
 \log X_N\leq\frac{\log C}{1-\sigma}
        +\sigma^j\log X_{K_0^jN},
\]
and the last term is bounded by
$\sigma^j(C'+s\log N+sj\log K_0)$, which tends to zero.  We conclude
that $X_N$ is uniformly bounded.

Equations \eqref{eq:sharp-nonzero-mode} and
\eqref{eq:sharp-zero-mode} now give
\eqref{eq:optimal-vertex-bound} for every $a\in S$ of weight at least
two.  To pass from cutoff estimates to a general finite-energy vector,
use $N=\max(1,\lceil h\rceil)$ on each input level $W(h)$ and add
squared norms: a fixed mode has different output energies for
different input energies.  Any specified Hermitian quasi-primary can
be included in such an $S$ by enlarging the fixed integer $D$.

For a weight-one generator, boundedness of $X_N$ gives
$R_a(N)\leq CN^{1/2}$.  For positive modes this is already a
half-order energy estimate.  For a negative mode $-m$, adjunction on
the output level gives a bound by $C(1+h+m)^{1/2}$ on $W(h)$, and
\[
 (1+h+m)^{1/2}\leq(1+m)^{1/2}(1+h)^{1/2}.
\]
Thus weight-one fields satisfy, in particular, at most linear energy
bounds.  Linear combinations of the Hermitian generators give the
same conclusion for all of $V_1$.  The Virasoro estimate was proved
in \eqref{eq:virasoro-linear-energy}.  This completes the proof of
Theorem~\ref{thm:energy}.

\begin{corollary}\label{cor:locality}
Every simple unitary strongly rational VOA $V$ is CKLW-strongly local.
Its smeared vertex operators generate an irreducible conformal net
$A_V$.
\end{corollary}

\begin{proof}
The unitary Virasoro decomposition expresses $V$ as a sum of
highest-weight representations.  Equivalently, at each energy the
orthogonal complement of the negative-Virasoro-mode descendants
consists of primary vectors.  These primary spaces are spanned by
Hermitian vectors. Hence the conformal vector and Hermitian primary
vectors generate $V$; see \cite[Remark~3.9]{CTW}. A finite family
suffices: decompose the vectors in $S$ into Virasoro descendants and
take Hermitian bases of the finitely many primary spaces that occur.
Together with $\nu$, these vectors generate $S$ and hence $V$.

Theorem~\ref{thm:energy} gives the optimal order $d-1$ for every such
primary of weight $d\geq2$, and supplies the required current and
Virasoro bounds.  The corrected strong-locality criterion
\cite[Theorem~3.7]{CTW}, with its Hermitian hypothesis
\cite{CTWCorrection}, therefore applies.  Finally
\cite[Theorem~6.8]{CKLW} constructs the irreducible conformal net from
this simple unitary strongly local VOA.
\end{proof}

We have constructed $A_V$ independently of any fusion positivity or
module integration.  In the following section it is this vacuum net
that supplies the annulus operators.
